% Также можно использовать \Referat, как в оригинале
\begin{abstract}
\begin{center}
\textbf{к выпускной квалификационной работе магитсра на тему «Исследование и разработка фреймворка для микроконтроллера MIK32»}

\textbf{Шеберстова Дмитрия Алексеевича}

\textbf{Руководитель - Заева Маргарита Анатольевна, к.т.н., НИЯУ МИФИ, ИИКС, доцент}

\end{center}

    Пояснительная записка к выпускной квалификационной работе магистра (диссертации) \pageref{LastPage}\,стр.%
    \ifnum \totfig >0
    , \totfig~рис.%
    \fi
    \ifnum \tottab >0
    , \tottab~табл.%
    \fi
    %
    \ifnum \totbib >0
    , \totbib~источн.%
    \fi
    %
    \ifnum \totapp >0
    , \totapp~прил.%
    \else
    .%
    \fi

ВСТРОЕННЫЕ СИСТЕМЫ, ИМПОРТОЗАМЕЩЕНИЕ, МИКРОКОНТРОЛЛЕР, ОТЕЧЕСТВЕННЫЕ ТЕХНОЛОГИИ, ФРЕЙМВОРК, I2C, GPIO, HAL, RISC-V

Актуальность исследования обусловлена необходимостью создания отечественных решений в области микроэлектроники и разработки инструментов программирования для российских микроконтроллеров. Микроконтроллер MIK32, построенный на открытой архитектуре RISC-V, представляет собой платформу для импортозамещения в критически важных отраслях.

Объектом исследования является микроконтроллер MIK32. Предметом исследования выступают методы и средства разработки программного обеспечения для данного микроконтроллера, включая архитектурные решения фреймворков и подсистемы управления периферийными устройствами.

Цель работы заключается в упрощении процесса разработки встроенных приложений за счет создания прототипа фреймворка для микроконтроллера MIK32, который обеспечит абстракцию аппаратных ресурсов.

Для достижения поставленной цели решались задачи: анализ архитектурных особенностей микроконтроллера MIK32 и существующих решений для RISC-V платформ; проектирование многоуровневой архитектуры фреймворка с четким разделением аппаратно-зависимых и аппаратно-независимых компонентов; разработка унифицированных подсистем управления GPIO и I2C; реализация драйверов для внешних расширителей портов; практическая апробация фреймворка в промышленных проектах.

Методы исследования включают: анализ архитектуры микроконтроллеров и существующих фреймворков; проектирование с применением паттернов абстракции; экспериментальное тестирование на реальных аппаратных платформах.

Научная новизна заключается в разработке архитектуры фреймворка, адаптированной к особенностям отечественного микроконтроллера MIK32, обеспечивающую совместимость с широким спектром периферийных устройств.

Практическая значимость подтверждается успешным применением фреймворка в двух промышленных проектах: системе управления нагревом Heating Unit и системе управления вентиляторами Midplane. Разработанный фреймворк способствует ускорению разработки встроенных систем на базе MIK32 и повышению конкурентоспособности отечественных технологий.

\end{abstract}

%%% Local Variables: 
%%% mode: latex
%%% TeX-master: "rpz"
%%% End: 
